\chapter{حدود الإجماع والقياس، وأدب المناظرة والاختلاف}

\section{التحذير من ادعاء الإجماع في مسائل الخلاف}
الحمد لله والصلاة والسلام على رسول الله، أما بعد؛ فقد استكمل الإمام الشافعي رحمه الله تعالى حديثه في مناقشة مخالفيه، محذراً من التساهل في ادعاء "الإجماع" في مسائل يوجد فيها خلاف معتبر بين السلف.

يقول الشافعي راداً على من يدعي الإجماع لنصرة مذهبه: «لست أقول -ولا أحد من أهل العلم- هذا مجمع عليه، إلا لما لا تلقى عالماً أبداً إلا قاله لك، وحكاه عمن قبله؛ كالظهر أربع، وكتحريم الخمر، وما أشبه هذا. وقد أجده يقول: "المجمع عليه"، وأجد من أهل العلم كثيراً يقولون بخلافه! وأجد عامة أهل البلدان على خلاف ما يقول: المجمع عليه!». 
فالإجماع اليقيني هو الإجماع على ما عُلم من الدين بالضرورة، وينبغي الحذر الشديد من ادعاء الإجماع الباطل لتمرير الآراء.

\section{قواعد في ضمان العاقلة لجناية الخطأ}
ناظر الشافعي رحمه الله أحد فقهاء الأحناف في مسألة جناية الخطأ وما تتحمله "العاقلة" (وهم عصبة الرجل وقرابته).
فالقاتل خطأً تتحمل عاقلته دية المقتول (مائة من الإبل)، لكن هل تتحمل العاقلة الجراحات الصغيرة التي لا تبلغ ثلث الدية (كالموضحة، وما دونها)؟
قال المناظر: لا تعقل العاقلة ما دون الموضحة؛ لأن رسول الله ﷺ لم يقض فيما دون الموضحة بشيء.
رد الشافعي بإلزام دقيق: إذا لم يقض فيها بشيء، فهل أهدرها؟ لا، بل الجراح مضمونة. وإذا كانت العاقلة تغرم الأكثر (الدية الكاملة)، فما المانع أن تغرم الأقل؟ فالقاعدة المطردة أن جميع ما كان جناية خطأ فعلى العاقلة، ولو كان درهماً يسيراً، طالما أن السبب هو "الخطأ". 
أما ما لزم الإنسان في غير جناية الخطأ (كالزكاة، والمهر، والكفارة، وجناية العمد، وما أفسدته مواشيه)، فكل هذا في ماله الخاص ولا تتحمله العاقلة أبدًا.

\section{مناظرات تطبيقية في القياس}

\subsection{دية العبد المقتول خطأً}
اختلفوا في جناية الحر على العبد خطأً؛ فقال الأحناف وأهل المدينة: "جراح العبد في ثمنه كجراح الحر في ديته" (أي تُحسب كنسبة مئوية). وخالفهم الشافعي، قائلاً إن جراح العبد تُقدر بما نقص من ثمنه (كأي سلعة تُتلف).
طالب المخالفُ الشافعيَّ بالدليل، ثم استدل المخالف بقول سعيد بن المسيب.
فرد عليه الشافعي بإنصاف وذكاء: أنت في أصولك لا ترى أقوال التابعين حجة، فكيف تحتج بها هنا؟ ثم أقام الشافعي القياس الصحيح مبيناً أن العبد يُفارق الحر في أن دية الحر مؤقتة شرعاً، بينما العبد ديته ثمنه، فأشبه السلع والبهائم المضمونة بقيمتها. وساق إلزاماً عقلياً وشرعياً للمخالف حتى أثبت تناقضه في قياسه.

\subsection{جواز استقراض الحيوان ورده بأحسن منه}
أنكر المخالف أن تُشترى الإبل بصفة في الذمة إلى أجل (كأن يستقرض بعيراً). فرد الشافعي بحديث أبي رافع رضي الله عنه: «أن النبي ﷺ استسلف من رجل بعيراً، فجاءته إبل فصرفه أن يقضيه إياه، فلم يجد إلا جملاً خياراً رباعياً، فقال ﷺ: أعطه إياه، فإن خيار الناس أحسنهم قضاءً». فبطل إنكارهم بالسنة الصريحة الجلية.

\section{الرخص والاستثناءات لا يُقاس عليها}
من درر القواعد الأصولية التي أرساها الشافعي رحمه الله: أن ما استُثني من أصل محرم للحاجة (الرخصة)، لا يجوز أن يُجعل أصلاً يُقاس عليه غيره.

\textbf{المثال الأول: المسح على الخفين:}
أصل فرض الله في الوضوء غسل القدمين، فلما مسح النبي ﷺ على الخفين، أخذناه اتباعاً ورخصةً لمن لبسهما على طهارة. ولم يجز لنا أن نقيس على الخفين غيرهما؛ فلا يجوز المسح على العمامة أو البرقع أو القفازين قياساً على الخفين، بل نبقي الأصل (الغسل) فيما عدا المستثنى بالنص.

\textbf{المثال الثاني: بيع العرايا:}
نهى النبي ﷺ عن بيع الثمر بالتمر (المزابنة) لوجود التفاضل والجهالة وكونه ربا. ولكنه رخص لبعض الفقراء الذين لا يملكون نقوداً أن يشتروا الرطب على النخل بخرصه من التمر ليأكلوه (في حدود خمسة أوسق)، وهذه هي "العرايا".
فالنهي عام ومستقر، والرخصة في العرايا خاصة، فلا يجوز أن نضرب الحديثين ببعضهما، ولا أن نقيس بيوع غرر أخرى على رخصة العرايا.

\section{أحكام تعبدية وأخرى معقولة المعنى}

\subsection{دية الجنين (الغرة)}
قضى النبي ﷺ في الجنين إذا أسقطته أمه ميتةً بجناية بغرة (عبد أو أمة، وتُقوّم بخمس من الإبل).
قال الشافعي: هذا حكم فارق حكم النفوس الأحياء؛ فلو سقط حياً ثم مات، ففيه دية كاملة (مائة إبل للذكر، وخمسون للأنثى). ولما حكم النبي ﷺ بالغرة للجنين الميت ولم يسأل أذكر هو أم أنثى؟ دل على أنه حكم "تعبدي"، سنه النبي ﷺ لحكمة يعلمها الله، فنسلم له، ولا نقيس عليه غيره. 
أما الأحكام المعقولة المعنى (التي ندرك علتها)، فنطيع الله فيها ونقيس عليها ما اشترك معها في ذات العلة.

\subsection{قاعدة الخراج بالضمان وحديث المصراة}
قضى النبي ﷺ أن «الخراج بالضمان»؛ فمن اشترى عبداً فاستغله وربح منه، ثم وجد فيه عيباً فرده، فإن المشتري يأخذ الربح ولا يرده للبائع؛ لأنه كان ضامناً للعبد لو مات. فقاس الشافعي على ذلك ثمر النخل وولد الجارية ولبن الماشية المشتراة.

لكنه استثنى "المصراة" (الشاة التي يُربط ضرعها ليجتمع اللبن فيغتر المشتري)، فإن النبي ﷺ أمر بردها ورد "صاع من تمر" معها كبدل للبن. فهنا نقف عند النص ولا نقيسه، ونعمل بالقياس العام في غير المصراة.

\section{أنواع الاختلاف (المحرم والسائغ)}
قال السائل: إني أجد أهل العلم قديماً وحديثاً مختلفين، فهل يسعهم هذا الاختلاف؟
قال الشافعي: الاختلاف من وجهين:
\begin{enumerate}
    \item \textbf{اختلاف محرم:} وهو ما أقام الله به الحجة في كتابه، أو على لسان نبيه ﷺ نصاً بيناً لا يحتمل التأويل؛ فهذا لا يحل الاختلاف فيه لمن علمه. (مثل من يعارض النصوص القطعية بالآراء العقلانية المحدثة).
    \item \textbf{اختلاف سائغ:} ما كان يحتمل التأويل ويُدرك قياساً، فذهب المتأول أو القايس إلى معنى يحتمله الخبر. فهذا لا أضيق فيه على المخالف كما أضيق في المنصوص، وهو من باب الاجتهاد المأجور صاحبه.
\end{enumerate}
وقد ذم الله التفرق بعد مجيء البينات ﴿وَمَا تَفَرَّقَ الَّذِينَ أُوتُوا الْكِتَابَ إِلَّا مِن بَعْدِ مَا جَاءَتْهُمُ الْبَيِّنَةُ﴾.

\section{معنى القرء في اعتداد المطلقة}
اختلف الصحابة رضي الله عنهم في تفسير قول الله تعالى: ﴿وَالْمُطَلَّقَاتُ يَتَرَبَّصْنَ بِأَنفُسِهِنَّ ثَلَاثَةَ قُرُوءٍ﴾. 
فقالت عائشة وزيد بن ثابت وابن عمر: الأقراء هي "الأطهار". 
وقال أبو بكر وعمر وعلي وابن عباس: الأقراء هي "الحيض".

وبيّن الشافعي أن اللفظ لغةً يحتمل الأمرين، لأن "القرء" يعني "الوقت" ويطلق عليهما معاً. لكن الشافعي رجح أن المراد به (الطهر)، مستدلاً باللغة؛ فمصدر (قرأ) بمعنى (جمع وحبس)، والرحم في حال الطهر يجمع الدم ويحبسه، وفي حال الحيض يرخيه ويرسله، فالطهر أولى باسم القرء.
كما استدل بالسنة؛ حيث أمر النبي ﷺ ابن عمر أن يطلق امرأته طاهراً من غير جماع، وقال: «فتلك العدة التي أمر الله أن يُطلق لها النساء».
وهذا مثال على الاختلاف السائغ الذي لا يُعنف فيه المخالف لتوفر أدلة محتملة للطرفين.